% GENERATED FILE. Edit canonical lessons, then run npm run generate:books.
% canonical-source-hash: fnv1a64:b9e65d0f1df217cd
% canonical-lessons: UR-C539-zabardasti, UR-C539-mukhtar, UR-C539-naguzir, UR-C539-salah, UR-C539-jhukana

\chapter{Free to Choose}
\label{ch:ur-free-to-choose}
\hlchaptermodality{539}

\begin{chapteropening}
\textbf{By the end of this chapter:} \emph{I can talk about free choices and being made to do things.}

Complete the last lesson of chapter 539: I can talk about free choices and being made to do things.
\end{chapteropening}

\section[\texorpdfstring{\ur{زبردستی} (zabardastī)}{زبردستی (zabardastī)}]{\ur{زبردستی} (zabardastī) — force (making someone do something)}

\label{lesson:UR-C539-zabardasti}



\begin{quote}
Before the new one: say the Urdu for to totter, then the Urdu for unshakable.
\end{quote}

\subsection*{You'll want to know: \ur{زبردستی}}
\textbf{\ur{زبردستی}} — \emph{zabardastī} — \textquotedblleft{}force (making someone do something)\textquotedblright{}.

\begin{quote}
\ur{اگر} \ur{وہ} \ur{نہیں} \ur{چاہتا} \ur{تو} \ur{زبردستی} \ur{ٹھیک} \ur{نہیں۔}
\end{quote}

\begin{quote}
\emph{agar voh nahīṅ chāhtā to zabardastī ṭhīk nahīṅ.}
\end{quote}

\begin{quote}
If he doesn't want to, forcing him isn't right.
\end{quote}

Feminine.

\begin{cousinweb}
\textbf{\ur{زبردستی}} (\emph{zabardastī}) is built on \emph{zabardast}, \textquotedblleft{}having the upper hand\textquotedblright{}, usually traced to Persian \emph{zabardast}, from \emph{zabar}, \textquotedblleft{}above\textquotedblright{}, and \emph{dast}, \textquotedblleft{}a hand\textquotedblright{}, with the noun ending \emph{-ī}.
\end{cousinweb}

\begin{grammarlens}[title={What you've built}]
One more word for saying what happens if.
\end{grammarlens}

\subsection*{Your turn}
\begin{itemize}
\raggedright
  \item \emph{Say it:} \emph{zabardastī}
  \item \emph{Say it:} \emph{zabardastī}, once more
  \item \emph{Say it:} \emph{agar voh nahīṅ chāhtā to zabardastī ṭhīk nahīṅ.} — If he doesn't want to, forcing him isn't right.
  \item \emph{Recall:} say the Urdu for to totter, then the Urdu for unshakable, then say \emph{zabardastī} again
\end{itemize}

\begin{tcolorbox}[breakable,colback=teal!4,colframe=teal!35!black,title={Before you move on}]
What does \ur{زبردستی} mean? (\textquotedblleft{}Force (making someone do something)\textquotedblright{}.) Say it once more.
\end{tcolorbox}
\section[\texorpdfstring{\ur{مختار} (mukhtār)}{مختار (mukhtār)}]{\ur{مختار} (mukhtār) — free to decide; empowered}

\label{lesson:UR-C539-mukhtar}



\begin{quote}
Before the new one: say the Urdu for unshakable, then the Urdu for force.
\end{quote}

\subsection*{You'll want to know: \ur{مختار}}
\textbf{\ur{مختار}} — \emph{mukhtār} — \textquotedblleft{}free to decide; empowered\textquotedblright{}.

\begin{quote}
\ur{اگر} \ur{آپ} \ur{مختار} \ur{ہیں} \ur{تو} \ur{آپ} \ur{کو} \ur{کون} \ur{روکے} \ur{گا؟}
\end{quote}

\begin{quote}
\emph{agar āp mukhtār haiṅ to āp ko kaun roke gā?}
\end{quote}

\begin{quote}
If the choice is yours, who will stop you?
\end{quote}

\begin{cousinweb}
\textbf{\ur{مختار}} (\emph{mukhtār}) is usually traced, through Persian, to Arabic \emph{mukhtār}, \textquotedblleft{}chosen; free to choose\textquotedblright{}.
\end{cousinweb}

\begin{grammarlens}[title={What you've built}]
One more word for saying what happens if.
\end{grammarlens}

\subsection*{Your turn}
\begin{itemize}
\raggedright
  \item \emph{Say it:} \emph{mukhtār}
  \item \emph{Say it:} \emph{mukhtār}, once more
  \item \emph{Say it:} \emph{agar āp mukhtār haiṅ to āp ko kaun roke gā?} — If the choice is yours, who will stop you?
  \item \emph{Recall:} say the Urdu for unshakable, then the Urdu for force, then say \emph{mukhtār} again
\end{itemize}

\begin{tcolorbox}[breakable,colback=teal!4,colframe=teal!35!black,title={Before you move on}]
What does \ur{مختار} mean? (\textquotedblleft{}Free to decide; empowered\textquotedblright{}.) Say it once more.
\end{tcolorbox}
\section[\texorpdfstring{\ur{ناگزیر} (nāguzīr)}{ناگزیر (nāguzīr)}]{\ur{ناگزیر} (nāguzīr) — unavoidable}

\label{lesson:UR-C539-naguzir}



\begin{quote}
Before the new one: say the Urdu for force, then the Urdu for free to decide.
\end{quote}

\subsection*{You'll want to know: \ur{ناگزیر}}
\textbf{\ur{ناگزیر}} — \emph{nāguzīr} — \textquotedblleft{}unavoidable\textquotedblright{}.

\begin{quote}
\ur{اگر} \ur{سفر} \ur{ناگزیر} \ur{ہے} \ur{تو} \ur{صبح} \ur{جلدی} \ur{نکلیے۔}
\end{quote}

\begin{quote}
\emph{agar safar nāguzīr hai to subah jaldī nikliye.}
\end{quote}

\begin{quote}
If the journey can't be avoided, set off early in the morning.
\end{quote}

\begin{cousinweb}
\textbf{\ur{ناگزیر}} (\emph{nāguzīr}) is usually traced to Persian \emph{nāguzīr}, \textquotedblleft{}with no way out\textquotedblright{}, built from \emph{nā}, \textquotedblleft{}not\textquotedblright{}, and \emph{guzīr}, \textquotedblleft{}escape\textquotedblright{}.
\end{cousinweb}

\begin{grammarlens}[title={What you've built}]
One more word for saying what happens if.
\end{grammarlens}

\subsection*{Your turn}
\begin{itemize}
\raggedright
  \item \emph{Say it:} \emph{nāguzīr}
  \item \emph{Say it:} \emph{nāguzīr}, once more
  \item \emph{Say it:} \emph{agar safar nāguzīr hai to subah jaldī nikliye.} — If the journey can't be avoided, set off early in the morning.
  \item \emph{Recall:} say the Urdu for force, then the Urdu for free to decide, then say \emph{nāguzīr} again
\end{itemize}

\begin{tcolorbox}[breakable,colback=teal!4,colframe=teal!35!black,title={Before you move on}]
What does \ur{ناگزیر} mean? (\textquotedblleft{}Unavoidable\textquotedblright{}.) Say it once more.
\end{tcolorbox}
\section[\texorpdfstring{\ur{صلاح} (salāh)}{صلاح (salāh)}]{\ur{صلاح} (salāh) — advice}

\label{lesson:UR-C539-salah}



\begin{quote}
Before the new one: say the Urdu for free to decide, then the Urdu for unavoidable.
\end{quote}

\subsection*{You'll want to know: \ur{صلاح}}
\textbf{\ur{صلاح}} — \emph{salāh} — \textquotedblleft{}advice\textquotedblright{}.

\begin{quote}
\ur{اگر} \ur{آپ} \ur{میری} \ur{صلاح} \ur{سنتے} \ur{تو} \ur{یہ} \ur{مشکل} \ur{نہ} \ur{ہوتی۔}
\end{quote}

\begin{quote}
\emph{agar āp merī salāh sunte to yih mushkil na hotī.}
\end{quote}

\begin{quote}
If you had listened to my advice, this trouble wouldn't have happened.
\end{quote}

Feminine.

\begin{cousinweb}
\textbf{\ur{صلاح}} (\emph{salāh}) is usually traced, through Persian, to Arabic \emph{ṣalāḥ}, \textquotedblleft{}soundness, what is right\textquotedblright{}; the sense \textquotedblleft{}advice\textquotedblright{} came later, and the sources we use do not show how it arose.
\end{cousinweb}

\begin{grammarlens}[title={What you've built}]
One more word for saying what happens if.
\end{grammarlens}

\subsection*{Your turn}
\begin{itemize}
\raggedright
  \item \emph{Say it:} \emph{salāh}
  \item \emph{Say it:} \emph{salāh}, once more
  \item \emph{Say it:} \emph{agar āp merī salāh sunte to yih mushkil na hotī.} — If you had listened to my advice, this trouble wouldn't have happened.
  \item \emph{Recall:} say the Urdu for free to decide, then the Urdu for unavoidable, then say \emph{salāh} again
\end{itemize}

\begin{tcolorbox}[breakable,colback=teal!4,colframe=teal!35!black,title={Before you move on}]
What does \ur{صلاح} mean? (\textquotedblleft{}Advice\textquotedblright{}.) Say it once more.
\end{tcolorbox}
\section[\texorpdfstring{\ur{جھکانا} (jhukānā)}{جھکانا (jhukānā)}]{\ur{جھکانا} (jhukānā) — to bow (one's head); to lower}

\label{lesson:UR-C539-jhukana}



\begin{quote}
Before the new one: say the Urdu for unavoidable, then the Urdu for a piece of advice.
\end{quote}

\subsection*{You'll want to know: \ur{جھکانا}}
\textbf{\ur{جھکانا}} — \emph{jhukānā} — \textquotedblleft{}to bow (one's head); to lower\textquotedblright{}.

\begin{quote}
\ur{اگر} \ur{وہ} \ur{شرمندہ} \ur{ہوتا} \ur{تو} \ur{سر} \ur{جھکاتا۔}
\end{quote}

\begin{quote}
\emph{agar voh sharmindā hotā to sar jhukātā.}
\end{quote}

\begin{quote}
If he were ashamed, he would hang his head.
\end{quote}

Past, with \ur{نے} (he or she …): \ur{اس} \ur{نے} \ur{سر} \ur{جھکایا}, \emph{us ne sar jhukāyā}.

\begin{cousinweb}
\textbf{\ur{جھکانا}} (\emph{jhukānā}) is the causative of \emph{jhuknā}, \textquotedblleft{}to bend\textquotedblright{}, which CDIAL lists under \emph{*jhukkati}, a form it reconstructs (CDIAL 5399).
\end{cousinweb}

\begin{grammarlens}[title={What you've built}]
One more word for saying what happens if.
\end{grammarlens}

\subsection*{Your turn}
\begin{itemize}
\raggedright
  \item \emph{Say it:} \emph{jhukānā}
  \item \emph{Say it:} \emph{jhukānā}, once more
  \item \emph{Say it:} \emph{agar voh sharmindā hotā to sar jhukātā.} — If he were ashamed, he would hang his head.
  \item \emph{Recall:} say the Urdu for unavoidable, then the Urdu for a piece of advice, then say \emph{jhukānā} again
\end{itemize}

\begin{tcolorbox}[breakable,colback=teal!4,colframe=teal!35!black,title={Before you move on}]
What does \ur{جھکانا} mean? (\textquotedblleft{}To bow (one's head); to lower\textquotedblright{}.) Say it once more.
\end{tcolorbox}
